\begin{abstract}
Python type errors depend on the state in which an operation executes.
A function call can change the type of an object attribute or a global
variable, making a previously valid operation incompatible with the updated
state. Detecting such errors requires connecting input conditions to the
state changes established across calls. We present an interprocedural
type-state analysis for Python type-error detection. Its central model
relates function entries to ordered operation requirements, state updates,
and exits, preserving both conditional effects and intermediate failures.
The analysis constructs these relations by sharing facts independent of
unknown entry types and using a large language model to propose implicit
type alternatives. At each call, it instantiates object identities and
propagates the corresponding state changes. A conflict triggers selective
expansion of the relevant callee type paths and contextual re-inference;
revised relations then update dependent callers. The analysis separately
reports runtime type errors and annotation-contract errors. We evaluate
paired detection on xxx recent GitHub bug-fix records, diagnostic precision,
component effects, and total analysis cost. The evaluation reports
xxx paired runtime detections and xxx paired annotation-contract detections,
with respective diagnostic precision of xxx\% and xxx\%.
\end{abstract}
